\section{DISCUSSÃO}

Os resultados indicam que a abordagem neuromórfica proposta é viável para predição de
parâmetros físico-químicos da água em ambiente amazônico, com desempenho consistente após
ajuste de hiperparâmetros e regularização por \textit{early stopping}. No cenário multimodal
sem vazamento de informação (sem uso das variáveis-alvo nas entradas), observou-se aumento da
dificuldade de predição, especialmente para turbidez, o que torna a avaliação mais realista
para uso em campo.

No benchmark com 10 sementes, Random Forest e LSTM apresentaram menores erros médios
globais, enquanto a SNN permaneceu competitiva e claramente superior ao MLP em robustez.
Os testes estatísticos reforçam essa leitura: não houve evidência de superioridade estatística
da SNN sobre Random Forest e LSTM para $\alpha=0{,}05$, mas houve vantagem significativa
da SNN sobre MLP. Portanto, os resultados sustentam a SNN como alternativa viável em
cenário de borda, ainda que não dominante em erro absoluto no conjunto atual.
Em termos espaciais, os mapas de pH, OD e turbidez revelam padrões coerentes de gradientes
e agrupamentos, o que reforça a consistência dos dados de campo e sugere que as variáveis
de entrada carregam informação relevante para a predição. As correlações entre temperatura,
cor aparente e turbidez observadas na análise exploratória estão alinhadas com processos
físico-químicos esperados em ambientes lênticos, fornecendo suporte adicional para o
modelo proposto.
Além do desempenho preditivo, a execução do modelo em Raspberry Pi 4B apresentou latência
sub-milissegundo por amostra, reforçando a viabilidade de inferência local em hardware de
baixo consumo e compatível com cenários remotos e offline.
\par
Do ponto de vista prático, a diferença entre modelos sugere que a escolha depende do objetivo
operacional. Random Forest e LSTM tendem a minimizar erro médio, enquanto a SNN oferece uma
alternativa robusta quando se prioriza inferência local com menor custo computacional e
potencial de eficiência energética. O comportamento mais instável do MLP confirma que modelos
densos clássicos são sensíveis em bases pequenas e podem variar de forma significativa entre
sementes. Além disso, a dificuldade sistemática na turbidez indica que esse parâmetro exige
informação adicional (por exemplo, variáveis hidrodinâmicas ou sazonais) para melhorar a
generalização.

\par
Outra consideração importante é que a codificação temporal por \textit{rate coding} implica
um compromisso entre fidelidade temporal e custo computacional: valores maiores de $T$ podem
aumentar a qualidade da representação, mas ampliam latência de inferência e consumo. Isso
sugere que, em cenários de uso real, a escolha de $T$ deve ser calibrada em função do contexto
operacional e das restrições energéticas do dispositivo de borda.
